\chapter{Evenredigheid}\label{ch:g7:prop}

In jaar 6 kwamen evenredigheidstabellen aan de orde (\cref{ch:g6:propdata}); in
dit hoofdstuk wordt \omterm{def:g7:prop:table}{evenredigheid} een volwaardig gereedschap: de kruisregel om een
vierde waarde te vinden, percentages, en schalen van kaarten en maquettes.
Snelheid, de beroemdste \omterm{def:g7:prop:table}{evenredigheid} van alle, komt in \cref{ch:g8:speed}.

\section{Tabellen herkennen en aanvullen}

\begin{definition}[Evenredigheidstabel]\label{def:g7:prop:table}
Een tabel met twee rijen is een
\emph{evenredigheidstabel}\index{evenredigheid} als de getallen van de tweede rij
die van de eerste zijn, vermenigvuldigd met één vast getal, de
\emph{constante}. \omterm{def:g3:numbers:evenodd}{Even} goed gezegd: alle kolomquotiënten
$\frac{\text{onder}}{\text{boven}}$ zijn gelijk.
\end{definition}

\begin{example}\label{ex:g7:prop:recognize}
\begin{center}
\renewcommand{\arraystretch}{1.3}
\begin{tabular}{l|ccc}
boven & $4$ & $10$ & $14$ \\
\hline
onder & $6$ & $15$ & $21$ \\
\end{tabular}
\end{center}
\omterm{def:g3:division:remainder}{Quotiënten}: $\frac64 = 1.5$, $\frac{15}{10} = 1.5$, $\frac{21}{14} = 1.5$:
\omterm{def:g6:propdata:def}{evenredig}, met constante $1.5$.
\end{example}

\begin{theorem}[Kruisregel]\label{thm:g7:prop:cross}
In een \omterm{def:g7:prop:table}{evenredigheidstabel} met kolommen $\begin{pmatrix} a \\ b \end{pmatrix}$ en
$\begin{pmatrix} c \\ d \end{pmatrix}$ zijn de kruisproducten gelijk:
\[
a \times d = b \times c .
\]
Daardoor kun je de vierde waarde uit de drie andere berekenen:
$d = \dfrac{b \times c}{a}$.
\end{theorem}

\begin{proof}
Zij $k$ de constante: $b = ka$ en $d = kc$. Dan is
\[
a \times d = a \times kc = k \times ac,
\qquad
b \times c = ka \times c = k \times ac :
\]
beide kruisproducten zijn gelijk aan $k \times ac$. $ad = bc$ door $a$ delen geeft
de formule voor $d$.
\end{proof}

\begin{example}[De vierde evenredige]\label{ex:g7:prop:fourth}
Als $7$ dezelfde boeken $2.8$ kg wegen, hoeveel wegen dan $12$ van die boeken?
Tabel en kruisregel:
\begin{center}
\renewcommand{\arraystretch}{1.3}
\begin{tabular}{l|cc}
boeken & $7$ & $12$ \\
\hline
kg & $2.8$ & $m$ \\
\end{tabular}
\end{center}
\[
7 \times m = 2.8 \times 12
\quad\Longrightarrow\quad
m = \frac{2.8 \times 12}{7} = \frac{33.6}{7} = 4.8 \text{ kg}.
\]
Controle via één stuk: één boek weegt $0.4$ kg, twaalf boeken wegen $4.8$ kg.
\end{example}

\section{Percentages}

\begin{method}[Percentages toepassen en vinden]\label{met:g7:prop:percent}
\begin{enumerate}
  \item $t\,\%$ \emph{toepassen}: vermenigvuldig met $\frac{t}{100}$ ($12\,\%$
        van $350$ is $0.12 \times 350 = 42$);
  \item het percentage \emph{vinden} dat een deel voorstelt: reken
        $\frac{\text{deel}}{\text{geheel}}$ uit en schrijf het op \omterm{def:g4:fractions:def}{noemer} $100$
        ($21$ van $60$: $\frac{21}{60} = \frac{35}{100} = 35\,\%$);
  \item vraag je altijd af: \emph{een percentage van wát?} Het geheel waarnaar je
        verwijst, doet ertoe.
\end{enumerate}
\end{method}

\begin{example}\label{ex:g7:prop:percent}
Op een school zijn $180$ van de $400$ leerlingen meisjes. Het percentage:
\[
\frac{180}{400} = \frac{180 \div 4}{400 \div 4} = \frac{45}{100}
= 45\,\% .
\]
En de andere kant op: hoeveel zijn de $55\,\%$ jongens?
$0.55 \times 400 = 220$. Controle: $180 + 220 = 400$.
\end{example}

\section{Schalen}

\begin{definition}[Schaal]\label{def:g7:prop:scale}
Op een kaart of maquette op \emph{schaal}\index{schaal} $\frac{1}{n}$ (geschreven
als $1 : n$) zijn de afstanden op de kaart \omterm{def:g6:propdata:def}{evenredig} met de echte afstanden, met
constante $\frac1n$:
\[
\text{afstand op de kaart} = \frac{1}{n} \times \text{echte afstand},
\]
en dat allebei gemeten in \emph{dezelfde eenheid}.
\end{definition}

\begin{example}\label{ex:g7:prop:scale}
Op een kaart $1 : 50\,000$ liggen twee steden $7$ cm van elkaar. De echte afstand:
\[
7 \times 50\,000 = 350\,000 \text{ cm} = 3\,500 \text{ m} = 3.5
\text{ km}.
\]
Andersom: een wandelpad van $12$ km meet op de kaart
$12$ km $= 1\,200\,000$ cm, dus $1\,200\,000 \div 50\,000 = 24$ cm.
\end{example}

\begin{omfigure}
\begin{tikzpicture}
\begin{axis}[omaxis,
  width=8.4cm, height=5.6cm,
  xmin=0, xmax=8.6, ymin=0, ymax=34,
  xlabel={kaart (cm)}, ylabel={echt (km)},
  xtick={2,4,6,8}, ytick={10,20,30}]
  \addplot[omDef, thick, domain=0:8.2] {4*x};
  \addplot[omDef, only marks, mark=*, mark size=1.5pt] coordinates
    {(2,8) (4,16) (6,24) (8,32)};
  \draw[dashed, omThm] (axis cs:7,0) -- (axis cs:7,28)
    -- (axis cs:0,28);
\end{axis}
\end{tikzpicture}

{\small Een \omterm{def:g7:prop:scale}{schaal} is een \omterm{def:g7:prop:table}{evenredigheid}: kaartafstanden tegen echte afstanden
liggen op één rechte \omterm{def:g6:lines:objects}{lijn} door de oorsprong (hier $1$ cm $\leftrightarrow$ $4$ km;
de stippellijnen lezen $7$ cm $\leftrightarrow$ $28$ km).}
\end{omfigure}

\begin{remark}
Grafieken zijn de snelste test op \omterm{def:g7:prop:table}{evenredigheid} (\cref{ch:g6:propdata}): punten op
één rechte \omterm{def:g6:lines:objects}{lijn} \emph{door de oorsprong} betekenen evenredige grootheden; een \omterm{def:g6:lines:objects}{lijn}
die de oorsprong mist (zoals een taxiprijs met een vast starttarief) betekent
\emph{niet} \omterm{def:g6:propdata:def}{evenredig}, ook al is het een rechte \omterm{def:g6:lines:objects}{lijn}. Affiene functies, in
\cref{ch:g9:linfunc}, maken dat precies.
\end{remark}

\section{Oefeningen}

\begin{exercise}[$\star$]\label{exo:g7:prop:1}
Welke tabellen zijn evenredigheidstabellen? Geef de constante als die bestaat.
\begin{center}
\renewcommand{\arraystretch}{1.2}
\begin{tabular}{l|ccc}
 & $3$ & $7$ & $11$ \\
\hline
 & $7.5$ & $17.5$ & $27.5$ \\
\end{tabular}
\qquad
\begin{tabular}{l|ccc}
 & $2$ & $5$ & $9$ \\
\hline
 & $6$ & $15$ & $28$ \\
\end{tabular}
\end{center}
\end{exercise}

\begin{exercise}[$\star$]\label{exo:g7:prop:2}
Vul aan met de kruisregel, en schrijf de berekening op: $5$ kg aardappelen kost $6$
euro; wat kost $8$ kg?
\end{exercise}

\begin{exercise}[$\star$]\label{exo:g7:prop:3}
Een printer print $36$ bladzijden in $3$ minuten. Hoeveel bladzijden print hij in
hetzelfde tempo in $5$ minuten? En hoe lang doet hij over $96$ bladzijden?
\end{exercise}

\begin{exercise}[$\star$]\label{exo:g7:prop:4}
Reken uit: $30\,\%$ van $250$; \quad $15\,\%$ van $60$; \quad de $t$ waarvoor
$t\,\%$ van $80$ gelijk is aan $20$.
\end{exercise}

\begin{exercise}[$\star$]\label{exo:g7:prop:5}
Van $25$ schoten scoorde een basketbalspeelster $18$ keer. Wat is haar
scoringspercentage? (Schrijf de \omterm{def:g6:fractions:def}{breuk} op \omterm{def:g4:fractions:def}{noemer} $100$.)
\end{exercise}

\begin{exercise}[$\star$]\label{exo:g7:prop:6}
Op een kaart $1 : 25\,000$ meet een pad $9$ cm. Wat is zijn echte lengte in km? Een
meer is $2$ km lang: hoe lang is het op de kaart?
\end{exercise}

\begin{exercise}[$\star$]\label{exo:g7:prop:7}
Een modelauto is gebouwd op \omterm{def:g7:prop:scale}{schaal} $1 : 43$. De echte auto is $4.3$ m lang. Hoe
lang is het model, in cm?
\end{exercise}

\begin{exercise}[$\star\star$]\label{exo:g7:prop:8}
Een recept gebruikt $240$ g chocolade voor $8$ porties. Aline heeft $300$ g
chocolade. Voor hoeveel porties is dat genoeg (een heel aantal porties)?
\end{exercise}

\begin{exercise}[$\star\star$]\label{exo:g7:prop:9}
De prijs van een jas daalt van $80$ naar $60$ euro.
\begin{enumerate}
  \item Hoeveel euro is de korting? En hoeveel procent van de oorspronkelijke
        prijs?
  \item Later gaat de prijs weer van $60$ naar $80$ euro. Leg uit waarom die
        stijging \emph{geen} $25\,\%$ is --- en reken het juiste
        stijgingspercentage uit.
\end{enumerate}
\end{exercise}

\begin{exercise}[$\star\star$]\label{exo:g7:prop:10}
Twee rijen van een tabel zijn \omterm{def:g6:propdata:def}{evenredig}:
\begin{center}
\renewcommand{\arraystretch}{1.2}
\begin{tabular}{l|ccc}
 & $4$ & $x$ & $10$ \\
\hline
 & $y$ & $10.5$ & $17.5$ \\
\end{tabular}
\end{center}
Zoek de constante, en daarna $x$ en $y$.
\end{exercise}

\begin{exercise}[$\star\star\star$]\label{exo:g7:prop:11}
Een kopieerapparaat verkleint documenten tot $80\,\%$ van hun maat. Een \omterm{def:g6:lines:objects}{lijnstuk} van
$10$ cm wordt gekopieerd, en daarna wordt de \emph{kopie} opnieuw gekopieerd met
dezelfde instelling.
\begin{enumerate}
  \item Hoe lang is het \omterm{def:g6:lines:objects}{lijnstuk} na één kopie? En na twee?
  \item Waarom is het antwoord na twee kopieën geen $60\,\%$ van het origineel?
        Welk enkel percentage hoort bij twee kopieën?
\end{enumerate}
\end{exercise}

\section{Opgave: de aarde meten met een stok}

\begin{problem}[{Weekendopgave --- schaduwen zijn een evenredigheid, en hoe
Eratosthenes in 240 v.Chr. de omtrek van de aarde uitrekende}]
\label{pb:g7:prop:1}
Tweeëntwintig eeuwen geleden mat de bibliothecaris van Alexandrië de aarde ---
zonder telescoop, zonder satelliet en zonder rekenmachine, met een stok, een put,
een kamelenkaravaan en de wiskunde van dit hoofdstuk. In deze opgave loop je zijn
stappen na: eerst de \omterm{def:g7:prop:table}{evenredigheid} van schaduwen, dan de beroemde berekening zelf,
en tot slot hoe zijn antwoord er op menselijke \omterm{def:g7:prop:scale}{schaal} uitziet.

\medskip
\noindent\textbf{Deel I --- De schaduw van een stok.}
De zon staat zo ver weg dat haar stralen ons \omterm{def:g6:lines:perp}{parallel} bereiken. Op een gegeven
moment zijn alle verticale voorwerpen en hun schaduwen dus \omterm{def:g6:propdata:def}{evenredig}.
\begin{enumerate}
  \item Om twaalf uur werpt een verticale stok van $1$ m een schaduw van $0.4$ m,
        en de schaduw van een boom meet $3.2$ m. Hoe hoog is de boom
        (\cref{thm:g7:prop:cross})?
  \item Op datzelfde moment werpt een toren een schaduw van $26$ m: hoe hoog is de
        toren? En hoe lang is de schaduw van een kind van $1.5$ m?
  \item Leg in één zin uit waarom al deze berekeningen alleen geldig zijn op
        \emph{hetzelfde moment} van de dag.
  \item De legende zegt dat Thales de Grote \omterm{def:g5:solids:def}{Piramide} mat door te wachten op het
        moment waarop \emph{zijn eigen schaduw precies zo lang was als hij zelf}.
        Wat is de evenredigheidsconstante op dat moment, en hoe hoog is de \omterm{def:g5:solids:def}{piramide}
        als de punt van haar schaduw op $147$ m ligt van het punt op de grond recht
        onder haar top?
  \item Vat deel I samen: welke grootheid is op één gegeven moment hetzelfde voor
        de stok, de boom, de toren en de \omterm{def:g5:solids:def}{piramide}
        (\cref{def:g7:prop:table})?
\end{enumerate}

\medskip
\noindent\textbf{Deel II --- De berekening van Eratosthenes.}
Eratosthenes kende twee feiten. In de stad Syene stond de zon om twaalf uur op de
dag van de zomerzonnewende \emph{precies in het zenit}: het zonlicht bereikte de
bodem van de diepste putten, en verticale stokken wierpen geen schaduw. In
Alexandrië, $800$ km recht naar het noorden, wierp een verticale stok op datzelfde
moment \emph{wel} een schaduw, en maakten de zonnestralen een hoek van $7.2^\circ$
met de verticaal.
\begin{enumerate}[resume]
  \item Leg uit waarom die twee waarnemingen samen bewijzen dat de grond van Syene
        en de grond van Alexandrië niet \omterm{def:g6:lines:perp}{parallel} zijn --- dat het oppervlak van de
        aarde dus gekromd is. (Wat zouden parallelle stralen met twee stokken op een
        platte aarde doen?)
  \item Op een tekening van de ronde aarde met parallelle zonnestralen duikt de
        $7.2^\circ$ van Alexandrië opnieuw op in het middelpunt van de aarde, als
        de hoek tussen de richtingen van de twee steden. Maak de tekening en
        overtuig jezelf (de nette verantwoording gebruikt de paren gelijke hoeken
        van \cref{ch:g7:triangles}).
  \item Welke \omterm{def:g6:fractions:def}{breuk} van een volle draai is $7.2^\circ$?
  \item De boog van Syene naar Alexandrië ($800$ km) hoort bij $7.2^\circ$; de hele
        \omterm{def:g6:measure:perimeter}{omtrek} hoort bij $360^\circ$. Zet de \omterm{def:g7:prop:table}{evenredigheid} op en reken de \omterm{def:g6:measure:perimeter}{omtrek} van
        de aarde uit.
  \item Leid daaruit de \omterm{def:g4:geometry:circle}{diameter} van de aarde af, met $\pi \approx 3.14$
        (\cref{prop:g6:measure:circle}) en afgerond op honderd kilometer. (De
        moderne waarde is $12\,742$ km --- hoe dichtbij kwam een man met een stok in
        240 v.Chr.?)
\end{enumerate}

\medskip
\noindent\textbf{Deel III --- De aarde op menselijke \omterm{def:g7:prop:scale}{schaal}.}
\begin{enumerate}[resume]
  \item Er wordt een globe gebouwd op \omterm{def:g7:prop:scale}{schaal} $1 : 40\,000\,000$. Laat met de \omterm{def:g6:measure:perimeter}{omtrek}
        uit vraag 9 zien dat de \omterm{def:g6:measure:perimeter}{omtrek} van de globe precies $1$ m is. Wat is haar
        \omterm{def:g4:geometry:circle}{diameter}, op de centimeter?
  \item De Mont Blanc rijst ongeveer $4.8$ km op. Zet $4.8$ km om in centimeters,
        deel door $40\,000\,000$, en geef de hoogte van de berg op de globe in
        millimeters. Wat zegt het antwoord over hoe ``hobbelig'' de aarde
        werkelijk is?
  \item Een wandelaar legt $40$ km per dag af. Hoeveel dagen doet hij in dat tempo
        over de hele \omterm{def:g6:measure:perimeter}{omtrek} van Eratosthenes --- en hoeveel jaar is dat ongeveer?
  \item De ademlucht van de atmosfeer zit vooral in de eerste $10$ km boven de
        grond. Welk percentage van de \omterm{def:g3:shapes:shapes}{radius} van de aarde (ongeveer $6\,370$ km) is
        dat (\cref{met:g7:prop:percent})? Rond af op een tiende procent.
  \item Historici schatten dat de waarde die Eratosthenes aankondigde, omgezet naar
        moderne eenheden, ongeveer $42\,000$ km kan zijn geweest. Neem $40\,000$ km
        als de echte waarde en reken zijn foutpercentage uit. Besluit met één zin
        over stokken, putten en \omterm{def:g7:prop:table}{evenredigheid}.
\end{enumerate}
\end{problem}
